\section{Where the Design Breaks}
\label{sec:discussion}

Every number in this section comes from one run of the command a user would run,
on the repository that contains \sgt{}, at a single commit. We took them together
rather than as we needed them, because the repository is still being worked on and
a set of counts gathered over a week is a set of counts of nothing. That repository holds
fifty-seven days and three hundred and forty-six commits of daily use by the
person who wrote the system. That is long enough to have found the limits below
and not long enough to have found the ones that need a year, and we say which is
which where we can tell. What emerges from the numbers is a distinction among
three properties a developer needs from any tool that rewrites files on their
behalf: that nothing is removed (the store is durable), that the tool operates on
what the developer named (the report is honest), and that new work can still be
recorded after the operation finishes.
Section~\ref{sec:discussion-synthesis} traces each failure class to the property
it violates; the numbers below are what those failures look like from inside a
session.

\subsection{The rate a developer meets}
\label{sec:session-rate}

A robustness harness exercised every write verb and every selector format across
289 randomised operation sequences totalling 10{,}237 applied operations. Of
those, 340 left the repository in a state that violated one of three invariants
(a file that does not rebuild from its records, a record that names a function no
file contains, or a composed file whose bytes differ from disk). The per-operation
violation rate is 3.3\%, but this number is the one we would choose if we wanted
the design to look good, because a developer does not run one operation. A
developer runs a session of twenty to fifty operations, and the question they
experience is whether that session contains at least one violation.

Of 289 sequences, 145 contain at least one violation targeting an entity the
developer would name---a per-session rate of 50\%. The violations cluster
rather than spreading evenly: 68\% of flagged operations target entities
(functions the developer named in the command), and of those entity-targeted
violations, 157 of 232 trace to a layout record being dragged along by the
entity operation rather than to the entity logic itself. The entity logic
worked; the seam between it and the surrounding whitespace did not. Splitting
the pool by target kind, operations aimed at entities the developer would name
(entity plus feature-level and filename-level selectors) carry a 3.1\% flag
rate; operations aimed at layout records alone carry 4.1\%. The seam is not
noisier than the core, but it produces a different kind of failure: a layout
violation leaves the developer's code correct and the tool's rebuild unable to
reproduce it, which is a durability failure that silently weakens the next
revert rather than a corruption the developer sees now.

The 50\% per-session rate sounds incompatible with the 3.3\% per-operation rate,
but the two numbers measure different properties and the gap between them is itself
a finding. Individual operations in isolation are reliable: a single revert fails at
1.1\%, a single save at 1.5\%, a single undo at 0.7\%. What fails is composition.
The round-trip probe---revert an operation, then restore everything the revert
removed, and check whether the bytes returned---fails at 44\% (175 of 399 trials).
A developer who runs one revert almost always gets the right files. A developer who
runs a revert, works for an hour, then decides they want it back discovers that the
restore does not return to the starting point nearly half the time. The mechanism is
that a revert's upward closure removes records whose dependents shift the composed
image, and a restore brings back the originals without the shifts, so the two
operations compose to a state neither one would produce alone. This is the
set-theoretic version of a well-known property in concurrent systems: individual
operations can be correct while their composition violates an invariant the
individual specifications do not constrain. The design claims that revert followed
by restore is an identity, and the harness measures how far from identity the
composition actually lands.

On the five real repositories (cloned from GitHub, never touched by \sgt{} before
the harness), four of five produced at least one violation, confirming that the
rate is a property of the design rather than an artefact of the study repository's
migration history. The one clean real repository (asyncer, 50 operations) was also
the smallest, so its survival is consistent with the hypothesis that violation
probability scales with store size rather than being a fixed defect rate.

The fixture arm and the real-repository arm measured different systems in a
precise sense. Fixtures are built to round-trip, so the write guard that fires
when the store does not reproduce the committed bytes---the guard responsible
for 70--79\% refusal on real repositories---was structurally unreachable in all
284 fixture sequences and fired 82 times in 137 real-repository operations.
Fixtures also write every caller and callee in one commit, so the
cross-commit dependency the \cmd{--keep-dependents} flag exists to preserve
was absent from all 878 fixture operations that tested it. A behaviour that
is absent from the corpus does not produce a zero rate---it produces no rate
at all, and a zero in the table reads as ``tested and clean'' when the honest
reading is ``not tested.'' The fixture arm is an excellent bug finder (nine
defect classes, all in the layout seam); it is not a deployment estimate, and
the 284-of-289 proportion in the pooled table would have hidden exactly that
distinction.

The practical consequence is blunt: a developer who uses the write layer daily
will encounter a violation within their first two sessions on average. Muir's
model of trust development explains why this timing is
catastrophic~\cite{muir1994}. Trust in automation passes through three
qualitatively different stages: predictability (can I anticipate what the system
will do?), dependability (does it do what it says it will?), and faith (will it
handle situations I have not yet seen?). A developer in their first two sessions
is still at the predictability stage---they are learning what the tool's
operations mean, not yet relying on them. A failure at this stage does not merely
subtract one positive experience from a base of many; it prevents the stage
transition entirely, because the developer cannot progress to depending on a
system whose behaviour they cannot yet predict. Lee and See's calibration model
explains why recovery is slow once trust is damaged~\cite{lee2004}; Muir's
developmental model explains why damage during formation is qualitatively
different from damage during maintenance---it is not a setback within a stage but
a barrier to reaching the next one. The timing is worse than it would have been
a year ago. Across the broader ecosystem, trust in AI-generated output is
actively declining---from 40\% to 29\% in one survey
cycle~\cite{stackoverflow2025}---while adoption is near universal. A developer
arriving at \sgt{} in 2026 already expects automation to be unreliable, which
means Muir's predictability stage requires more confirmations before the
transition than it did when the baseline expectation was neutrality rather than
scepticism. The threshold for earning trust rises as the ecosystem trains
developers to withhold it.

A concrete instance from this paper's own construction illustrates the pattern
the harness measures in the abstract. During development, a merge conflict
resolution took the in-progress copy of six files whole, and the copies predated
several fixes that had already landed. Git reported the merge as a success; the
working tree was clean; six functions were silently removed, including one that
the agent interface calls on every restore operation. Two of the six had no
failing test, because the same resolution had removed their tests along with
their code. A suite that passes over reduced coverage reads the same as a suite
that passes over full coverage, and a green checkmark after a merge is the same
silent-success pattern the pilot participant encountered at the interface level.
Finding the removal required comparing name-by-name coverage across thirty files
against a reference commit---exactly the operation \sgt{} is designed to
perform, and one we did not think to run on our own merge because we had no
reason to doubt the result.

\subsection{Most of a repository is not functions}
\label{sec:coverage}

\cmd{sgt log --summary} prints \cmd{396 file(s), 6821 symbol(s), 154 feature(s),
25\% entity coverage}, and we printed that last figure for months before checking
it against the records it claims to summarise. It divides the files that still hold
a live function-level record at the current commit by every path any record
mentions, including files deleted long ago, so its two halves are drawn from
different populations and the ratio understates the coverage by a factor of nearly
three. The count below is the one to read.

Counting the records themselves says something different. This store holds records for
534 files, more than the 396 on disk today because a record outlives the file it was
read from, and of those 534, 381 carry function-level records and 153 are recorded
whole. The whole-file set is 82 Markdown files, 19 JSON, 15 \LaTeX{}, 13 HTML, and a
scattering of CSS, SVG, and configuration, and it contains no Python and no
TypeScript, so every file in a language \sgt{} parses has a function-level record. The
design's coverage of code is complete, and the files left over, nearly three in ten of
everything the store has a record for, hold no function for anything to be filed
under.

So the limit is the repository's composition. Calling it a parser problem let us describe a fixed property of software as though it were a defect we could fix. Go and Rust have
functions. A tree-sitter grammar plus a decision about which node kinds count as
one is a day of work per language, and we would take both tomorrow. It would not
move that share at all, because there is no Go and no Rust in this repository. What is in it is
prose, configuration, and generated files, in a proportion no repository of this kind
escapes. A change to four lines of a deployment manifest has no function to file the
developer's request under, no matter how many grammars we ship. Every operation in Section~\ref{sec:design}
still runs on those files, at git's grain, and none of them gets sharper than
\cmd{git revert} would be. A reader deciding whether to adopt \sgt{} should ask how
much of the work they need to \emph{take apart} happens inside functions, which is a
different question from how much of their repository is code. In the repository
this paper studies, 70\% of files are code and 100\% of the tangling problems the
scenarios describe live in functions, because a deployment manifest does not
entangle with a retry policy inside a shared definition.

\subsection{Where a function loses its history}
\label{sec:identity}

The design assumes a function's record survives when its source is renamed or
moved, and that assumption rests on a matcher with a threshold in it. \sgt{} links a function across a commit first by
its name, then by identical bytes, then by identical structure with formatting and
comments ignored, and last by token overlap above 0.80, with a guard that rejects
any pair whose token counts differ by more than a factor of two. The 0.80 comes
from sem, which uses the same figure for the same job~\cite{sem2026}. It scores the
overlap between the two bodies' sets of distinct tokens, which puts the breaking
point lower than the number sounds: replacing more than about one token in nine drops
a pair below it, and the function is then recorded as one removed and a different one
added.

The developer then sees a function whose recorded history begins at that commit, with
its earlier edits still filed under a group the function is no longer part of, and an
agent asked to clean up a module produces exactly this, renaming six functions and
rewriting their bodies in one sitting. The detectors that get a function's origin right
compare full syntax
trees across the two versions~\cite{godfrey2005,kim2005,falleri2014,tsantalis2018},
where \sgt{} uses the token count because it runs on every save and has to finish while
the developer is waiting. Lowering the threshold would link functions that merely
resemble each other, and a wrong link is the worse failure, because it files one
function's edits under another function's name and gives the developer no reason to
doubt the answer. \cmd{sgt advanced identity join} states the link by hand, and using
it requires the developer to notice that a function's history restarted, which is the
kind of thing they installed \sgt{} in order not to have to notice.

\subsection{Fifty-six files we cannot rebuild}
\label{sec:unreproducible}

The same command reports 56 files whose bytes on disk no valid set of records can
produce. The denominator is the 381 files whose records name functions and not the
534 the store holds, because a file recorded whole holds its own bytes and cannot
fail to come back, so one function-level file in seven does not rebuild.
\sgt{} keeps every one of them, prints the count, names the first five, and points at
\cmd{sgt advanced fsck --tree}, so nothing was deleted and nothing was quietly wrong
(Section~\ref{sec:sets}). What the developer loses is the record for those files: a
revert of a group that touched one of them leaves the file as it is, and \cmd{sgt show}
on a function inside it answers from an incomplete chain.

We have not tracked down what put all 56 in that state, and since this store holds
records written by five different versions of our miner in fifty-seven days, some of
them are the residue of our own format changes. A developer adopting \sgt{} on a
repository with two years of history should expect a number of this kind on the first
run, and should read it as work the design has not finished.

A deeper instance of the same incompleteness surfaced during the evaluation. \sgt{}
reads a repository's history in ten-second chunks and records a single commit as
its stopping point. A single commit cannot describe position in a branching
history, so when a chunk boundary falls between a merge and its second parent,
commits on the side branch are stepped over and never read. On one repository
mined twice, the second reading was missing 49 operations from one commit, and
8 commits carrying roughly 30 files were absent from both readings. Both readings
reported themselves complete. The same mechanism means that a function renamed
across a chunk boundary is recorded as one name removed and a different name
added, producing 108 differing identifiers between two readings of the same
commit on another repository. This is a self-reporting failure of the same class as
the ones the pilot found at the interface level: the tool claims completeness while
the claim is false, and nothing in the tool's own output gives the developer
reason to doubt.

A deeper problem compounds the chunk-boundary effects: the budget is wall-clock
time, so the store's contents depend on machine load at the moment of mining.
Five mines of one repository produced 150, 168, 168, 237, and 237 operations from
the same commit, and the two large stores reconstructed fewer files than the small
ones, because extra operations introduced by a longer chunk sometimes compose to
the wrong bytes. The practical consequence is that two developers onboarding the
same repository on different machines can end up with different stores, different
reconstruction rates, and different sets of refused operations, with no indication
that anything distinguishes them. A design that records \emph{what happened} should
not depend on \emph{when it was read}, and this coupling is a bug rather than a
tradeoff.

The consequence on external repositories is severe. On the first repository
outside the author's own projects that the evaluation pointed \sgt{} at (746
commits, 118 merges, two package renames), reconstruction reproduces 35\% of
function-level files. This is not noise around a successful result; it is a total
failure of the write layer, because the guard that protects composition refuses
every write verb when the composed image disagrees with disk. A developer who
clones a mid-sized library, runs \cmd{sgt init}, and waits the six unattended
minutes it takes for onboarding to reach genesis gets a tool that can attribute
but cannot act---the read layer works, and every write verb refuses. The
reconstruction rate on the author's own repository (7 of 8 function-level files
rebuild) is not evidence that the system works; it is evidence that the author's
habits (linear history, no deletions, no package renames) happen to avoid the
three inputs that break composition.

The rebuild also fails in the opposite direction: on 28 of 29 corpus
repositories, the composed image materialises files that do not exist at
\cmd{HEAD}---1{,}864 paths across the corpus, produced by operations whose
targets were deleted or renamed and whose records were never pruned. A metric
that checks only whether a file at \cmd{HEAD} can be reproduced from its records
cannot see this failure, because the invented file is not at \cmd{HEAD} and
therefore not in the denominator. Requiring the rebuild to produce
\cmd{HEAD}'s file set \emph{and nothing more} moves the corpus median rebuild
rate from 0.33 to 0.24. The two numbers answer different questions: the first
asks how much of the working tree the record can explain, and the second asks
how close the record comes to reproducing the working tree as a whole. Both are
needed, because the write guard compares the entire composed image against
disk, and a surplus file disagrees just as a missing function does. On the five
real repositories, surplus paths account for a third of refusals---the write
verb refuses because the rebuild contains files that no longer exist, not
because it cannot reproduce the ones that do.

The onboarding path has a separate failure that compounds the completeness problem.
\sgt{} reads history in ten-second chunks, and between chunks nothing tells the
developer their history is partial: no command drives the walk to completion, the
incompleteness flag exists only in the JSON output humans do not read, and the
summary prints ``13 file(s) on disk differ from the recorded state---\cmd{sgt save}
absorbs them.'' Following that instruction would record 746 commits of other
people's work as a single save by whoever just typed it, destroying the provenance
the tool exists to produce. A confident instruction pointing at a destructive
action, firing on every repository large enough to matter, is a stronger instance
of the self-reporting pattern than any single defect the robustness harness found.
The refusal messages themselves misdirect: in two cases on the recovery path, the
tool names a wrong cause (``operation does not exist'' for one it holds; ``set
\cmd{OPENAI\_API\_KEY}'' for a failure unrelated to any credential). A system
whose failure messages are wrong has the honesty property the design claims for
it in the code paths that were tested, and not as a structural guarantee.

\subsection{The layout seam}
\label{sec:layout-seam}

Order and whitespace are encoded as symbols but are not symbols under the
kernel's relations. A blank line between two functions is recorded as a layout
record, carries no callers and no dependencies, participates in no structural
community, and cannot fork because nothing competes for it. Yet composition
treats it identically to a function body: including a feature includes its
layout records, reverting one removes them, and the rebuild-from-records
guard checks the composed bytes against disk with no distinction between a
wrong function body and a wrong separator. This encoding preserves insertion
order deterministically---without it, recomposed files would lose their
formatting---and it is the single most productive source of defects across
the entire evaluation.

The evidence arrives from three directions at once. In the robustness harness,
157 of 232 entity-targeted violations trace to a layout record being dragged
along by the entity operation rather than to any failure in the entity logic
itself. On real repositories at rest---before any operation runs---19, 48, and
89 orphaned layout records appear across three of five clones, growing to
hundreds after the mine completes, because the fork rule withdraws the entity
records that anchor them and leaves the layout records pointing at a gap.
And the evaluation's only two recoverability failures---a 33-byte file emptied
while the tool reported it had done nothing---both targeted layout records. No
entity record lost a byte in 10{,}237 operations. Every byte that was lost came
through the seam.

The mechanism is not a bug in any individual component. The kernel's relations
(calls, co-change, structural containment) ground entity records in a
dependency structure that composition can reason about: if a function depends
on another, including one includes the other. Layout records have no such
grounding. They float between the entities they separate, and when one of those
entities is withdrawn by the fork rule the layout record stays behind, orphaned,
pointing at a position that no longer exists in the composed image. The composed
file then disagrees with disk by the bytes of that orphan, the rebuild guard
fires, and every write verb that touches that file refuses. The developer sees
``file does not rebuild from its records'' and has no way to know that the record
in question is a blank line between two functions that the fork rule withdrew.

A designer facing this choice has three options. Record order as metadata
attached to each entity (loses the ability to revert a formatting change
independently). Record order as a relation between adjacent entities rather
than a symbol in its own right (the approach Darcs takes with its sequencing
theory, at the cost of a different composition algebra). Or accept the seam as a
maintenance cost and fix the orphans it produces. The third is what we have been
paying, and the evaluation's finding is that the cost is higher than we
estimated: not a cosmetic imprecision in the rebuild, but the mechanism behind
half the addressable failure surface and the only path to data loss the harness
found.

\subsection{Findable is not decomposable}
\label{sec:findable-not-decomposable}

On the two study repositories, every information request is answerable from the
\sgt{} record in at most three commands: eight derivability checks, run three
times across two rebuilds, pass every time. A reader might stop there. The checks
are too easy. They ask whether the answer is \emph{reachable}---whether some
sequence of \cmd{sgt log}, \cmd{sgt show}, and \cmd{sgt find} can surface the
relevant symbols---and at function granularity almost anything is reachable. The
operations the design is built on are not retrieval but \emph{decomposition}:
naming a coherent piece of work as one unit, and reverting that unit without side
effects. Measured as decomposition, the same trees fail. One feature in each
repository participates in 18 of 22 episodes (coursecraft's \texttt{CLI Scaffold})
or 20 of 22 (confplan's \texttt{Schedule Grid}), so reverting that feature removes
work from nearly every sitting. Five of 34 features carry a label that describes
under a third of their symbols. A 5-symbol removal reverses 13 edits, because the
structural cut crosses the episode boundary wherever a hub function participates.

The three numbers are not three findings. Every episode splitting across 3--6
features, one feature spanning 90\% of episodes, and reverting 5 symbols costing
13 edits are one fact seen from three sides: the clustering cuts along the code
(files, coupling) and the work ran along features (sessions, requests), so the cuts
land crosswise. A paper that reports them as a list of weaknesses would pad the
count and hide that a single design choice produces all of them. The choice is
commitment~1: the function as the unit of recording. A function belongs to one
feature at a time. A feature is any set of functions. When a hub function
participates in every episode, every episode and every feature overlap in that hub,
and the overlap is the mechanism behind every number above.

The corpus evaluation confirms the same mismatch at a larger scale. On four
real repositories (two external, two this project's own transcripts), one
developer request becomes a median of 4--6 \sgt{} features and up to 12 on the
longest session. No repository reaches 1:1. The finding kills a framing the paper
once used: features are not recovered intents. At this granularity they are
recovered \emph{sub}-intents---structural communities that correspond to what a
developer did at the code level, not to what they asked for at the request level.
Suchman's distinction between plans and situated actions explains why no
refinement of the clustering will close this gap~\cite{suchman1987}: a
developer's request is a plan in her sense, a resource that oriented the agent's
work rather than a specification that determined it, and one orienting statement
produces several independent implementation decisions the statement never
mentioned. The gap between request and feature is the gap between planning and
doing.
The distinction matters because the paper's scenarios ask ``which request produced
this function,'' and the feature map answers a different question: ``which other
functions were modified in the same structural neighbourhood.'' That answer is
useful---it is the cross-file grouping git cannot provide---but it is a weaker
claim than intent recovery, and the paper must state it as the claim actually
supported rather than letting the reader infer the stronger one.

The practical consequence for the developer in Section~\ref{sec:scenarios} is
that \cmd{sgt revert "price cache"} does not do what the scenario promises. The
price cache was one request. The tool split it into three features: one holding
the cache helper, one holding the TTL logic because it co-changed with an
unrelated timer, and one holding the parameter that threads the TTL through
\sym{db.py::query} because that function's coupling graph places it in the
database cluster. The developer who typed one sentence must now name three
features to remove one request, and deciding which three requires reading the
feature map---the same comprehension labour the tool was supposed to eliminate.
The labour is smaller (read a feature tree instead of reading the code), and it
is faster (a \cmd{sgt log --map} call instead of a \cmd{git log -L} call per
function), but it is not zero, and a design whose selling point is ``name it and
remove it'' must say honestly that ``name it'' sometimes means ``find the three
pieces it became.'' Clark's grounding theory explains the gap: shared vocabulary
requires proposal and acceptance between two
parties~\cite{clark1996}, and the tool proposes its structural
categories without any acceptance step from the developer. The merge command
exists to repair this, but a developer who does not know the pieces exist cannot
know they need merging. Surfacing the mismatch---showing, at save time, which
features the tool proposes and asking the developer to confirm---is the fix this
design has not yet built. Building it requires interrupting the save flow,
which is the forcing-function cost the fork rule already pays. That fix also
contradicts Section~\ref{sec:names}'s own premise: Shipman and Marshall showed
that systems requiring users to formalise their work get
abandoned~\cite{shipman1999}, and a confirmation prompt at every save is a
formalisation request at the highest-frequency operation. The design space is
therefore narrower than Clark's model alone suggests: the acceptance step has to
be cheap enough that developers do not skip it, and intermittent enough that it
does not become the same cognitive tax the tool was built to eliminate. A preview
that appears only when the grouping changed since the last save, or a deferred
confirmation that accumulates until the developer acts, might thread that needle.
We have not built either, and until we do the grounding gap stays open.

\subsection{What a colleague sees}
\label{sec:colleague}

Because the records are files in the repository, everyone who pulls the repository
gets them whether they asked for them or not, and this section works through what
that means for three people who are not the developer running \sgt{}: a colleague
who has never installed it, a colleague who installed a different version of it,
and the reviewer of a pull request.

The one we worried about most turned out to be the easy case. A colleague who has
never heard of \sgt{} commits ordinary code, and the next \cmd{sgt sync} mines
their commits into function-level records, the same way \sgt{} mines pre-existing
history. Their edits arrive at the right grain and group with the work they belong
to. What does not arrive is the request sentence, because there was none to
capture: the name for their group is drawn from commit subjects instead. A
colleague who writes \cmd{fix tests} and \cmd{wip} gets names derived from the
code, the same failure Section~\ref{sec:names} describes for a save with no
message. The grain \sgt{} recovers from anybody; the vocabulary only from somebody
who was running it.

The harder case is a colleague who does install it and runs a different version.
\sgt{} recovers its records from code instead of being handed them, so the program
that did the reading is part of what a record is, and an edit's identifier is
computed from its content together with the version of the miner that produced it.
Two versions therefore mint different identifiers for the same edit, and a sync
across that gap would quietly alias records that mean different things, which \sgt{}
avoids by refusing the whole union and printing what to do. Two people using git
have to agree about bytes and nothing else, and two people using \sgt{} have to
agree about a program. That refusal
is correct and it is also the cost we are least able to speak to, because
Section~\ref{sec:unreproducible} reports five miner versions inside fifty-seven days
and every one of those transitions would have stopped a collaborator. It stopped
nobody, since there was nobody to stop, and a design decision that a solo repository
charges nothing for is a decision that repository cannot evaluate.

The third case is the reviewer, and here we have no answer at all. A pull request that
carries \sgt{} records shows them as what they are on disk, a few hundred JSON files
whose names are hashes, and no reviewer reads that; they review the code, as they
should, and how anybody should read a large model-written change spanning many files is
an open question being studied on its own~\cite{heander2026}. The reason is structural,
not incidental. Code review evaluates correctness: does this code do what it should?
The \sgt{} records encode provenance: which request produced which function? Evaluating
provenance requires the reviewer to read the original requests, compare them against the
grouping, and judge whether the inference is right---which is the comprehension work
the tool was built to eliminate, performed by a person who did not write the code and
has even less context than the developer who did. A reviewer who cannot evaluate the
record cannot catch a wrong resolution or a mislabelled feature, so the record that
makes \cmd{revert} work is not part of what anybody approved, and a resolution one
developer reasoned through and a label another corrected both enter the shared branch
unexamined. The rule in Section~\ref{sec:fork} asks a person to settle a fork, and on
a team that person is still one person, while the review that would have caught them
being wrong reads a diff their decision does not appear in.

\subsection{The fork rule has never met two people}
\label{sec:forkoff}

Section~\ref{sec:fork} says a resolved fork has to pass the project's test command
before \sgt{} will include it, and the summary above also prints
\cmd{oracle: unconfigured}. In the repository we use every day, resolving a fork
means the developer picks a version and \sgt{} takes their word for it, because
nobody ever pointed \sgt{} at our test runner. A safety property a developer has to
configure is a safety property most repositories will not have, so the strongest
claim in Section~\ref{sec:fork} holds only where somebody did that setup, and the
person who wrote the rule did not do it in the repository where the rule runs every
day. The fix needs no research, because \sgt{} could find the test command the
project already uses and offer it during \cmd{sgt init}, and making the developer
ask for it was the wrong default.

The condition the rule has been tested under is narrower still. This repository's
history carries three git identities and one human being behind all of them, so
every merge \sgt{} has ever seen had the same person on both sides of it, which is
the condition in which the rule is most obviously right, because the person a fork
stops is the person who wrote both versions and can answer it in seconds. It is
also the condition in which the rule is cheapest. On a repository where twenty
people edit the same functions every day, the same rule produces a queue of
decisions whose two sides were written by people who have not spoken, and Nelson et
al.'s account of how teams work through merge conflicts describes the barriers a
queue of that kind meets~\cite{nelson2019}. We expect that team would switch the
rule off and take the textual merge, and our evidence cannot tell them not to.

The setting we would build applies git's rule at the grain of a function, merging
when the two edits changed no overlapping lines and asking when they did, which
would also let \sgt{} report a divergence while both sides are still being written
rather than when they meet~\cite{guimaraes2012}. We have not built it, because we
cannot tell a reader which of the two should be the default, and the repository we
would have tested it on has one person in it. This is the part of the design a
study of larger teams would most usefully overturn.

\subsection{The record the agent already had}

The grouping \sgt{} recovers by clustering existed before any of the code did. The
agent held it while it worked, and every agent we have used discards it when the
session ends. Hutchins's account of distributed cognition explains what
disappeared~\cite{hutchins1995}: the knowledge of which functions serve which
purpose was formerly distributed across three supports---the developer's
procedural memory of typing the code, the commit structure that recorded
boundaries, and the naming conventions that made intent readable from identifiers.
Agent-mediated development removes the first two supports (the developer typed
sentences, not code; the commit happened once at the end, not at each boundary)
and leaves the third intact but insufficient, since function names describe
implementation rather than purpose. \sgt{} reconstructs what was lost by mining
the structural channels (coupling, co-change) that survive the transition.
A record written by the agent, one entry per request naming the functions it
changed, would make the inference in Section~\ref{sec:names} unnecessary, and it
would fix the case in Section~\ref{sec:walkthrough} where a rename nobody asked
for got a name derived from the code.

Half of that record we do take, and the half we take is the half a vendor cannot
spoil. Where the developer's agent is Claude Code, \sgt{} installs a hook on
prompt submission and stores every sentence the developer typed, with its time.
The words that named the work are on disk before any code exists, and \cmd{sgt now}
answers with what was asked for rather than with a count of edits. That is twenty
lines and one vendor's hook interface, and the same twenty lines have to be
written again for the next agent, which is the price we decided we could pay,
because what they capture is the developer's own English and a change of vendor
cannot make yesterday's sentence wrong.

The half we do not take is the agent's account of which functions each request
touched, and the reason is that a record of that kind is a claim about the files
which we would have to either trust or check. Trusting it puts the correctness of
a revert inside a transcript nobody validated against the code. Tang et
al.'s finding that inaccurate self-reporting accounts for 22.6\% of all
developer--agent misalignment episodes---rising even as other failure classes
decline~\cite{tang2026agentfailures}---says directly how often that trust would
be misplaced. Checking it means
reading the files and computing the answer, at which point we did not need the
account. What the agent's account would improve is the label, which is the part of
\sgt{} that is allowed to be wrong, and improving a label is worth a hint and not
worth a dependency. So the split we have arrived at is that membership is read
from code and names are read from words, and the words come from the hook where
there is a hook and from the save message where there is not. What we do not know
is whether an agent's account of its own edits is accurate enough to move
membership too, and that is the measurement we would ask a reader of this paper to
make before any other.

\subsection*{Open science}

\sgt{} is open source, and the repository it versions is the repository that
implements it, so the artefact and the evidence in this section are the same
object: a reader can clone it, check out the commit the artefact names, run
\cmd{sgt log --summary}, and read every number in this section off the output. Run
at a later commit the counts will have moved, since the repository is the working
one. The within-subjects study in Section~\ref{sec:eval-study} is pre-registered: its
design, measures, and the thresholds in Table~\ref{tab:measures} were fixed before
recruiting, so that the first row remains the first row after we have seen the
data. Two moderated pilots (Section~\ref{sec:pilot}) validated the tasks, the
rubric, and the scoring script; they found twelve tool defects and four design
defects, all fixed before the full study ran. The two repositories the study
uses, the transcripts that produced them, and the correctness script are
released with the results. The study's recorded data are screen activity,
task outcomes, and interview audio from consenting professional developers, with no
personal data beyond what consent and compensation require.
